\clearpage
\appendix

\section{Software and Computational Environment} \label{sec:a-Software}

All experiments were conducted in Python. The key software dependencies used
in this replication are listed in Table~\ref{tab:software}. The complete
environment specification, including all transitive dependencies, is provided in the accompanying code repository.

\begin{table}[h]
\centering
\caption{Software dependencies used in this replication.}
\label{tab:software}
\begin{tabular}{lll}
\hline
\textbf{Package} & \textbf{Version} & \textbf{Purpose} \\
\hline
Python      & 3.10.10 & Programming language \\
NumPy       & 2.2.6 & Numerical computation \\
PyTorch     & 2.12.0 & Neural network implementation \\
Matplotlib  & 3.10.9 & Visualisation \\
SciPy       & 1.15.3 & Signal processing and interpolation \\
tqdm        & 4.67.3 & Training progress tracking \\
\hline
\end{tabular}
\end{table}

\paragraph{Hardware.}
Experiments were run on a machine equipped with an Intel(R) Core(TM) Ultra 7 265 (20 cores) processor (dedicated CPU usage oscillated between 95\% and 100\%), 32 GB of RAM,
and Intel(R) Graphics graphics card. Training of the RL agent required approximately
55 hours for 700'000 episodes.

\section{Controller performances}\label{sec:a-controll_perf}

\subsection{Computation of Percentage Gaps}

To compare controllers across indicators with different units and magnitudes,
relative percentage gaps were computed with respect to the best-performing
controller for each demand scenario and performance indicator. The convention
was chosen such that \(0.0\%\) always denotes the best observed result, while
positive values indicate worse performance relative to that best result.

For indicators where higher values are better, the percentage gap was computed
as:

\begin{equation}
    \mathrm{Gap}
    =
    \frac{x_{\mathrm{best}} - x}{|x_{\mathrm{best}}|}
    \cdot 100
\end{equation}

where \(x\) is the value obtained by a given controller and
\(x_{\mathrm{best}}\) is the highest value observed among all controllers for
the same demand scenario.

For indicators where lower values are better, the percentage gap was computed
as:

\begin{equation}
    \mathrm{Gap}
    =
    \frac{x - x_{\mathrm{best}}}{|x_{\mathrm{best}}|}
    \cdot 100
\end{equation}

where \(x_{\mathrm{best}}\) is the lowest value observed among all controllers
for the same demand scenario.

The higher-is-better indicators are \(\mathrm{VKT}\), \(\bar{v}\),
\(\bar{v}_{\mathrm{queue}}\), and episode reward. The lower-is-better indicators
are \(\mathrm{VHT}\) and \(\mathrm{VHT}_{\mathrm{queue}}\).

For the episode reward, larger values are better because the reward is defined
as a negative penalty on deviations from the desired control-cell density. Thus,
a value closer to zero represents better density tracking. Results are available in Table \ref{tab:rankings}.

\clearpage

\subsection{Full controller performance results for each indicator and noise demand scenario
}
\begingroup
\scriptsize
\setlength{\tabcolsep}{3pt}
\renewcommand{\arraystretch}{1.12}
\begin{longtable}{p{5.2cm}rrrrrrr}
\caption{Full controller performance results across base and noisy demand scenarios.}\label{tab:full_results}\\
\toprule
Controller & \makecell{Noise std\\veh/h} & \makecell{VKT\\veh$\cdot$km} & \makecell{VHT\\veh$\cdot$h} & \makecell{VHT$_{\mathrm{queue}}$\\veh$\cdot$h} & \makecell{$\bar{v}$\\km/h} & \makecell{$\bar{v}_{\mathrm{queue}}$\\km/h} & \makecell{Episode\\reward} \\
\midrule
\endfirsthead
\caption[]{Full controller performance results across base and noisy demand scenarios. (continued)}\\
\toprule
\textbf{Controller} & \makecell{\textbf{Noise std}\\\textbf{veh/h}} & \makecell{\textbf{VKT}\\\textbf{veh$\cdot$km}} & \makecell{\textbf{VHT}\\\textbf{veh$\cdot$h}} & \makecell{\textbf{VHT$_{\mathrm{queue}}$}\\\textbf{veh$\cdot$h}} & \makecell{\textbf{$\bar{v}$}\\\textbf{km/h}} & \makecell{\textbf{$\bar{v}_{\mathrm{queue}}$}\\\textbf{km/h}} & \makecell{\textbf{Episode}\\\textbf{reward}} \\
\midrule
\endhead
\midrule
\multicolumn{8}{r}{Continued on next page}\\
\endfoot
\bottomrule
\endlastfoot
No control & 0 & 27700.62 & 309.79 & 309.79 & 89.42 & 89.42 & -6852.17 \\
ALINEA ($K_R$=10) & 0 & 27700.62 & 232.49 & 415.94 & 119.15 & 66.60 & -2281.23 \\
ALINEA opt $\bar{v}$ ($K_R$=23) & 0 & 27700.62 & 232.05 & 381.98 & 119.37 & 72.52 & -2211.65 \\
ALINEA opt $\bar{v}_{\mathrm{queue}}$ ($K_R$=21) & 0 & 27700.62 & 232.07 & 381.00 & 119.36 & 72.70 & -2213.92 \\
PI-ALINEA ($K_R$=100, $K_P$=4) & 0 & 27700.62 & 233.74 & 419.41 & 118.51 & 66.05 & -2481.82 \\
PI-ALINEA opt $\bar{v}$ ($K_R$=2, $K_P$=93) & 0 & 27700.62 & 231.12 & 404.59 & 119.85 & 68.47 & -2062.83 \\
PI-ALINEA opt $\bar{v}_{\mathrm{queue}}$ ($K_R$=2, $K_P$=56) & 0 & 27700.62 & 231.74 & 351.09 & 119.53 & 78.90 & -2162.09 \\
RL controller (ep 263366) & 0 & 27696.70 & 231.09 & 465.37 & 119.85 & 59.52 & -2062.83 \\
\midrule
No control & 50 & 27690.87 & 309.06 & 309.06 & 89.60 & 89.60 & -6854.04 \\
ALINEA ($K_R$=10) & 50 & 27690.87 & 232.42 & 415.22 & 119.14 & 66.69 & -2284.73 \\
ALINEA opt $\bar{v}$ ($K_R$=23) & 50 & 27690.87 & 232.03 & 385.96 & 119.34 & 71.75 & -2222.98 \\
ALINEA opt $\bar{v}_{\mathrm{queue}}$ ($K_R$=21) & 50 & 27690.87 & 232.04 & 385.56 & 119.34 & 71.82 & -2224.51 \\
PI-ALINEA ($K_R$=100, $K_P$=4) & 50 & 27690.87 & 233.59 & 418.68 & 118.55 & 66.14 & -2471.97 \\
PI-ALINEA opt $\bar{v}$ ($K_R$=2, $K_P$=93) & 50 & 27690.87 & 231.04 & 422.34 & 119.85 & 65.57 & -2064.03 \\
PI-ALINEA opt $\bar{v}_{\mathrm{queue}}$ ($K_R$=2, $K_P$=56) & 50 & 27690.87 & 231.66 & 352.03 & 119.53 & 78.66 & -2163.82 \\
RL controller (ep 263366) & 50 & 27684.80 & 230.98 & 473.81 & 119.86 & 58.43 & -2065.24 \\
\midrule
No control & 100 & 27619.04 & 308.53 & 308.53 & 89.52 & 89.52 & -6861.00 \\
ALINEA ($K_R$=10) & 100 & 27619.04 & 231.74 & 418.67 & 119.18 & 65.97 & -2280.96 \\
ALINEA opt $\bar{v}$ ($K_R$=23) & 100 & 27619.04 & 231.42 & 396.43 & 119.34 & 69.67 & -2230.15 \\
ALINEA opt $\bar{v}_{\mathrm{queue}}$ ($K_R$=21) & 100 & 27619.04 & 231.42 & 392.87 & 119.35 & 70.30 & -2229.82 \\
PI-ALINEA ($K_R$=100, $K_P$=4) & 100 & 27619.04 & 233.00 & 420.11 & 118.54 & 65.74 & -2481.98 \\
PI-ALINEA opt $\bar{v}$ ($K_R$=2, $K_P$=93) & 100 & 27619.04 & 230.47 & 418.16 & 119.84 & 66.05 & -2077.39 \\
PI-ALINEA opt $\bar{v}_{\mathrm{queue}}$ ($K_R$=2, $K_P$=56) & 100 & 27619.04 & 230.95 & 355.10 & 119.59 & 77.78 & -2154.37 \\
RL controller (ep 263366) & 100 & 27612.90 & 230.51 & 459.77 & 119.79 & 60.06 & -2093.03 \\
\midrule
No control & 150 & 27573.66 & 305.89 & 305.89 & 90.14 & 90.14 & -6831.61 \\
ALINEA ($K_R$=10) & 150 & 27573.66 & 231.39 & 407.10 & 119.17 & 67.73 & -2291.89 \\
ALINEA opt $\bar{v}$ ($K_R$=23) & 150 & 27573.66 & 231.16 & 395.99 & 119.28 & 69.63 & -2255.83 \\
ALINEA opt $\bar{v}_{\mathrm{queue}}$ ($K_R$=21) & 150 & 27573.66 & 231.15 & 391.24 & 119.29 & 70.48 & -2253.44 \\
PI-ALINEA ($K_R$=100, $K_P$=4) & 150 & 27573.66 & 232.92 & 411.97 & 118.38 & 66.93 & -2536.98 \\
PI-ALINEA opt $\bar{v}$ ($K_R$=2, $K_P$=93) & 150 & 27573.66 & 230.12 & 408.89 & 119.82 & 67.44 & -2089.26 \\
PI-ALINEA opt $\bar{v}_{\mathrm{queue}}$ ($K_R$=2, $K_P$=56) & 150 & 27573.66 & 230.80 & 361.12 & 119.47 & 76.36 & -2198.59 \\
RL controller (ep 263366) & 150 & 27564.70 & 230.20 & 434.14 & 119.74 & 63.49 & -2114.65 \\
\midrule
No control & 200 & 27673.15 & 304.38 & 304.38 & 90.92 & 90.92 & -6799.73 \\
ALINEA ($K_R$=10) & 200 & 27673.15 & 232.30 & 410.42 & 119.13 & 67.43 & -2292.58 \\
ALINEA opt $\bar{v}$ ($K_R$=23) & 200 & 27673.15 & 232.23 & 425.02 & 119.16 & 65.11 & -2281.76 \\
ALINEA opt $\bar{v}_{\mathrm{queue}}$ ($K_R$=21) & 200 & 27673.15 & 232.19 & 422.52 & 119.18 & 65.49 & -2274.67 \\
PI-ALINEA ($K_R$=100, $K_P$=4) & 200 & 27673.15 & 233.08 & 406.37 & 118.73 & 68.10 & -2418.12 \\
PI-ALINEA opt $\bar{v}$ ($K_R$=2, $K_P$=93) & 200 & 27673.15 & 230.97 & 455.91 & 119.81 & 60.70 & -2079.41 \\
PI-ALINEA opt $\bar{v}_{\mathrm{queue}}$ ($K_R$=2, $K_P$=56) & 200 & 27673.15 & 231.56 & 377.14 & 119.51 & 73.38 & -2173.60 \\
RL controller (ep 263366) & 200 & 27668.50 & 231.08 & 436.15 & 119.74 & 63.44 & -2102.57 \\
\midrule
No control & 250 & 27630.56 & 311.39 & 311.39 & 88.73 & 88.73 & -7037.18 \\
ALINEA ($K_R$=10) & 250 & 27630.56 & 232.32 & 432.74 & 118.94 & 63.85 & -2355.20 \\
ALINEA opt $\bar{v}$ ($K_R$=23) & 250 & 27630.56 & 232.00 & 412.64 & 119.10 & 66.96 & -2304.48 \\
ALINEA opt $\bar{v}_{\mathrm{queue}}$ ($K_R$=21) & 250 & 27630.56 & 231.99 & 407.66 & 119.10 & 67.78 & -2302.28 \\
PI-ALINEA ($K_R$=100, $K_P$=4) & 250 & 27630.56 & 233.10 & 425.50 & 118.53 & 64.94 & -2481.05 \\
PI-ALINEA opt $\bar{v}$ ($K_R$=2, $K_P$=93) & 250 & 27630.56 & 230.71 & 461.15 & 119.76 & 59.92 & -2097.64 \\
PI-ALINEA opt $\bar{v}_{\mathrm{queue}}$ ($K_R$=2, $K_P$=56) & 250 & 27630.56 & 231.27 & 399.37 & 119.48 & 69.19 & -2187.24 \\
RL controller (ep 263366) & 250 & 27622.60 & 230.80 & 449.73 & 119.68 & 61.42 & -2125.13 \\
\end{longtable}
\endgroup


